% Level 2, Week 8 -- Testing: Prove It Works (answer key)
\answersweek{8}{Testing: Prove It Works}{Answers for workbook chapter 8 and homework 8.}

\begin{teachernote}[How the test-writing checks work]
In Exercises 8.1, 8.2 and Homework 8, \texttt{starter.py} is the student's \emph{test file}. \texttt{check\_my\_work.py} runs the student's
tests against the correct function (they must all pass) and against three deliberately broken versions written inside \texttt{test\_\ldots.py}
(each broken version must make at least one of the student's tests fail). Any tests that do this are correct: the solutions below are one example.
\end{teachernote}

\section{Workbook exercises}
\begin{answer}[Exercise 8.1 Test the Clamp]
\textbf{Step 1:} 3 tests pass; the empty ones pass because a test with no \texttt{assert} checks nothing, so it cannot fail.
\textbf{Step 3:} with four tests, the last line looks like \texttt{1 failed, 3 passed in 0.03s}. \textbf{Broken clamps:} \enquote{forgets the low limit} is caught by any
too-small test, \enquote{forgets the high limit} by any too-big test, \enquote{swaps the limits} by either.
\end{answer}
\pyfile{../exercises/week08/ex-test-the-clamp/solution.py}

\begin{answer}[Exercise 8.2 Test the Best Score]
\textbf{Planning table:} \texttt{best\_score([])} expects \texttt{0} (edge); a list with the best last, e.g. \texttt{[5, 1, 8]} expects \texttt{8} (edge);
\texttt{[40]} expects \texttt{40} (edge); \texttt{[7, 7, 7]} expects \texttt{7} (edge). \textbf{Which test catches which bug:} the normal test catches
\enquote{returns the first score}; the empty-list test catches the crash; only a best-is-last test catches \enquote{skips the last score}.
\end{answer}
\pyfile{../exercises/week08/ex-test-best-score/solution.py}

\begin{answer}[Exercise 8.3 Fix the Bugs the Tests Found]
Each fix is marked \texttt{FIX 1} to \texttt{FIX 4} in the comments. Accept \texttt{return max(0, lives - damage)} for the first one. Check that no student edited
\texttt{test\_fix\_the\_bugs.py} (the checks only count if the tests are unchanged).
\end{answer}
\pyfile{../exercises/week08/ex-fix-the-bugs/solution.py}

\section{Homework}
\begin{answer}[Core: Test the Cipher]
\texttt{encode("xyz", 3)} is \texttt{"abc"}. Step 3 example: \texttt{"go team!"} with shift 1 is \texttt{"hp ufbn!"}.
The wrap-around test catches the no-wrap bug; the space test catches \enquote{loses spaces}; any normal test catches \enquote{goes backwards}.
\end{answer}
\pyfile{../exercises/week08/hw-test-the-cipher/solution.py}

\begin{answer}[Challenge: Spot the edge]
Many answers are right. Examples: \texttt{count\_vowels("dragon")} is 2 (normal); edges: \texttt{""} is 0, \texttt{"rhythm"} is 0, \texttt{"AEIOU"} (capitals).
\texttt{lives\_left(3, 1)} is 2 (normal); edges: \texttt{lives\_left(1, 5)} is 0, \texttt{lives\_left(0, 0)} is 0.
\texttt{last\_item(["a", "b", "c"])} is \texttt{"c"} (normal); edges: a one-item list, an empty list (decide: crash or \texttt{None}?).
\end{answer}

\begin{answer}[Creative]
Look for: a file \texttt{my\_tests.py} next to the program, an import of the student's own function, two tests with \texttt{assert},
and an edge case that is really an edge (empty, zero, first/last, on a limit).
\end{answer}
